\chapter*{Lời Tựa}
\addcontentsline{toc}{chapter}{Lời Tựa}
\markboth{Lời Tựa}{Lời Tựa}

\begin{truyen}{Lời Tựa}{What We Understand Too Late}
\chuhoa{C}{ó những bài học không đến lúc ta còn thong thả, mà đến khi lòng đã đau, tay đã hụt, và thời gian đã đi quá xa để quay lại như cũ.}
\textit{(Some lessons do not arrive while life is calm, but only after pain has struck, something has slipped away, and time has gone too far to turn back unchanged.)}

Quyển hai mươi tám này chọn một chủ đề ai cũng từng chạm tới: những điều ta hiểu khi đã muộn. Muộn để giữ một người, muộn để chăm một thân thể, muộn để sửa một lời nói, muộn để kịp nắm một cơ hội.
\textit{(This twenty-eighth volume chooses a theme everyone has touched: what we understand too late. Too late to keep someone, too late to care for the body, too late to fix a sentence, too late to seize an opportunity.)}

Nhưng hiểu muộn không phải vô ích. Hiểu muộn có thể cứu phần đời còn lại khỏi đi lặp lại cùng một sai lầm. Vì thế, 50 câu chuyện trong quyển này không nhằm kể khổ, mà nhằm giúp người đọc dừng lại sớm hơn, nghĩ sâu hơn, và sống tử tế hơn.
\textit{(But late understanding is not useless. It can still save the rest of life from repeating the same mistakes. That is why the 50 stories here are not meant to complain about life, but to help readers pause sooner, think deeper, and live more kindly.)}
\end{truyen}

\begin{baihoc}
Điều đáng sợ nhất không phải là học muộn, mà là học rồi vẫn sống y như cũ.
\textit{(The most frightening thing is not learning late, but learning and still living exactly the same.)}
\end{baihoc}

\begin{ghinhoanh}
\textbf{Short English to remember:}

Late lessons can still save the rest of your life.

Understand sooner. Regret less.
\end{ghinhoanh}

\begin{tuhocnhanh}
1. Điều gì trong đời em đang sợ một ngày nào đó sẽ hiểu ra quá muộn? \textit{(What in your life are you afraid you may understand too late one day?)}

2. Nếu phải chọn một điều để thay đổi ngay sau khi đọc quyển này, em sẽ chọn điều gì? \textit{(If you had to change one thing right after reading this volume, what would it be?)}
\end{tuhocnhanh}

\begin{flushright}
\textbf{Tôi Nguyễn Văn Sang}\\
\textit{GV THPT Nguyễn Hữu Cảnh - TP HCM}
\end{flushright}

\ngancach